\documentclass[12pt]{article}
\usepackage[margin=1in]{geometry}
\usepackage[T1]{fontenc}
\usepackage{xcolor}
\usepackage{soul}
\definecolor{garnet}{HTML}{73000A}

% Highlight commands using soul (supports line breaks)
\newcommand{\ylw}[1]{{\sethlcolor{yellow}\hl{#1}}}
\newcommand{\rd}[1]{{\sethlcolor{red!30}\hl{#1}}}
\newcommand{\blu}[1]{{\sethlcolor{cyan!30}\hl{#1}}}

\setlength{\parindent}{0pt}
\setlength{\parskip}{0.3em}

\begin{document}

\noindent
{\large\textbf{REVIEW KEY:}}\\
\ylw{Yellow} = Filler (words without benefit) \quad
\rd{Red} = Unsupported claim \quad
\blu{Blue} = Too much technical detail, not enough about student

\vspace{0.5em}
\noindent
{\color{garnet}\rule{\textwidth}{0.5pt}}
\vspace{0.5em}

\noindent
Department of Mechanical Engineering\\
University of South Carolina\\
Columbia, SC 29208\\
\today

\vspace{0.5em}

\noindent
SC Space Grant Consortium\\
College of Charleston\\
Charleston, SC 29424

\vspace{0.5em}

\noindent
\textbf{Re: Faculty Recommendation for Noah Young --- Undergraduate Research Award}

\noindent
{\color{garnet}\rule{\textwidth}{0.5pt}}

\vspace{0.3em}

\noindent
Dear Members of the Evaluation Committee:

I recommend Noah Young for the SC Space Grant Undergraduate Research Award. As his faculty advisor in the Integrated Multiphysics and Systems Engineering Laboratory (iMSEL) since August 2025, I have supervised his research on radar simulation and autonomous perception systems from his first week of college. \rd{His output in five months has been exceptional for any undergraduate}; for a freshman, \ylw{it is remarkable}.

Noah joined iMSEL at the start of his first semester as an Honors College aerospace engineering student. He is now co-authoring an ASME IDETC 2026 paper on synthetic radar data augmentation, in which he developed a method using land/water segmentation and object extraction to regenerate radar imagery from real sensor data. \blu{This required him to understand radar signal geometry, terrain interaction effects, and image processing pipelines}---concepts he learned on the job while simultaneously taking his first college courses. He maintains a 4.0\,GPA.

His technical contributions span multiple systems in our lab. He fine-tuned a YOLO model for maritime small-object detection and deployed it on the NVIDIA Jetson Orin, \blu{optimizing model size and inference latency for embedded hardware with constrained memory and power}. He designed a 3D-printable radar antenna mount in Creo Parametric that enabled our Furuno NXT radar to be integrated onto a research vessel for offshore data collection---\blu{a mechanical design task that required understanding vibration loads, waterproofing constraints, and cable routing}. He built the raw radar-to-PNG conversion pipeline and real-time visualization tools that our entire team now uses for experimental work. He annotated IR/RGB image pairs for multi-spectral training datasets, \blu{a task that requires understanding both the data and the detection models that will consume it}.

One aspect of Noah's work that I want to emphasize: he delivers weekly technical presentations and monthly written reports to our Department of Defense sponsors. This is a requirement we place on all lab members, but \rd{Noah's reports stand out for their clarity}. \rd{His figures are publication-quality}, \rd{his writing is structured}, and \rd{he can explain technical results to audiences who are not specialists in his specific area}. \rd{This communication skill is uncommon at any level} and \ylw{indicates genuine understanding rather than surface familiarity}.

His proposed project---an open-source terrain-occluded radar simulator validated against measured returns---addresses a real infrastructure gap. \blu{High-fidelity electromagnetic solvers require expensive licenses and hours of computation per scenario. Simplified point-target models ignore terrain interactions critical for low-altitude systems.} Noah proposes a middle path: \blu{geometric ray-casting against SRTM/NASADEM digital elevation models, with configurable radar parameters matched to our X-band specifications, validated against 4--5 real radar scenarios from our offshore campaigns. The key deliverable is documented accuracy bounds: quantitative metrics establishing when geometric approximation suffices and when it fails.} \blu{I have reviewed his proposed validation framework, which adapts the double validation metric from automotive radar simulation to terrain-radar interactions. The approach is sound, the real data exists in our lab, and the C++ acceleration target of 10,000+ ray-terrain intersections per second is achievable with established spatial indexing techniques.}

\blu{The NASA connection is direct. Mars 2020 Perseverance achieved 5-meter landing accuracy using Terrain Relative Navigation---a system that was validated extensively in simulation before deployment. Lunar landers under the Artemis program face similar terrain-sensing challenges. Noah's simulator would provide accessible, validated infrastructure for developing and testing terrain-relative navigation algorithms without requiring expensive electromagnetic solvers or repeated field campaigns.}

Noah is an AIAA member and participates in the USC Fly Club, \ylw{reflecting genuine engagement with aerospace beyond coursework}. He is MathWorks-certified in Simscape and proficient in Python, C++, and MATLAB. \ylw{I am confident he will execute the proposed 400-hour research plan and deliver the stated open-source tool, validation dataset, and symposium presentation.}

\ylw{I recommend Noah Young for this award without hesitation.}

\vspace{0.8em}

\noindent
Sincerely,

\vspace{1.5em}

\noindent
Dr.\ Yi Wang\\
Associate Professor\\
Department of Mechanical Engineering\\
University of South Carolina\\
\textit{ywang@sc.edu}

\newpage

\section*{REVIEW SUMMARY}

\subsection*{\ylw{YELLOW} --- Filler / Padding}

\begin{enumerate}
\item \textbf{``it is remarkable'':} Generic superlative that adds nothing beyond ``exceptional'' already stated in the same sentence.
\item \textbf{``indicates genuine understanding rather than surface familiarity'':} Vague interpretation. The specific behaviors should speak for themselves.
\item \textbf{``reflecting genuine engagement with aerospace beyond coursework'':} Opinion gloss. Just state the memberships; let reviewers draw conclusions.
\item \textbf{``I am confident he will execute...'':} Boilerplate confidence statement. Redundant if previous paragraphs showed his track record.
\item \textbf{``I recommend Noah Young for this award without hesitation'':} Redundant with opening sentence. Standard closing filler.
\end{enumerate}

\subsection*{\rd{RED} --- Unsupported Claims}

\begin{enumerate}
\item \textbf{``His output in five months has been exceptional for any undergraduate'':} ``Exceptional'' needs quantification. Compared to what baseline? How many undergrads have you supervised?
\item \textbf{``Noah's reports stand out for their clarity'':} Stand out compared to whom? Show, don't tell: quote a sponsor's comment or describe what makes one figure clear.
\item \textbf{``His figures are publication-quality'':} Which figures? Has one actually been published? ``Publication-quality'' is meaningless without context.
\item \textbf{``his writing is structured'':} Structured how? What approach does he use that others don't?
\item \textbf{``he can explain technical results to audiences who are not specialists'':} Give an example. Did a non-technical sponsor praise his presentation?
\item \textbf{``This communication skill is uncommon at any level'':} Huge claim with zero evidence. Uncommon among whom---freshmen? Graduate students? Faculty? Provide a reference frame.
\end{enumerate}

\subsection*{\blu{BLUE} --- Too Much Technical Detail}

\begin{enumerate}
\item \textbf{Para 2, second half:} ``radar signal geometry, terrain interaction effects, and image processing pipelines'' focuses on concepts, not student growth. What did he struggle with? How did he learn?
\item \textbf{Para 3, task descriptions:} Lists technical tasks (YOLO tuning, Jetson deployment, antenna mount, pipeline) but describes SYSTEMS more than the STUDENT. How did he approach debugging? Did he fail first? Work independently?
\item \textbf{Para 5, nearly all:} 7 sentences about the simulator's technical approach (ray-casting, SRTM/NASADEM, validation metrics, C++ performance). Explains the PROJECT, not why Noah is qualified or what he'll learn. Reviewers fund students, not software.
\item \textbf{Para 6, entirely:} NASA mission examples (Perseverance, Artemis) explain why the tool is useful but say nothing about Noah. This is project justification that belongs in his proposal, not character evidence in your letter.
\end{enumerate}

\subsection*{OVERALL ASSESSMENT}

\textbf{Strengths:} Specific timeline (August 2025), concrete deliverable (ASME co-authorship), quantitative achievement (4.0 GPA), sponsor interaction detail (weekly/monthly to DoD), named certification (MathWorks Simscape).

\textbf{Critical weaknesses:}
\begin{itemize}
\item Paras 5--6 are entirely about technical approach and NASA relevance---should describe what makes NOAH capable.
\item Communication skills paragraph is all tell, no show. Add ONE concrete example.
\item No failure/growth narrative. Where did he struggle? What did he learn from mistakes?
\item Superlatives without comparison frames need context (e.g., ``one of three freshmen in 15 years I've trusted with sponsor-facing presentations'').
\end{itemize}

\textbf{Fix:} Replace paras 5--6 with student characterization. Add one specific anecdote to the communication paragraph. Reframe para 3 tasks to show student qualities, not system specs.

\end{document}
